\section{Rough figure inventory for collaboration review}

This appendix is compiled only in the internal rough-build variant. Its purpose is to make the expected visual footprint of the paper visible before every scientific gate is closed. Figures in this appendix are therefore candidates, diagnostics or interpretation aids rather than additional publication-authorising claims. Numerical statements remain governed by the main chapters and their stated evidence boundaries.

\subsection{Additional gated analysis views}

\GatedFigure{figures/gated/fig08_mc_deltaE_E_I.png}{0.92}{Preliminary Sample~I MC $\Delta E$--$E$ view retained for layout and analysis review.}{Quarantined analysis output; useful for seeing the intended comparison, not for numerical citation.}

\GatedFigure{figures/gated/fig08_mc_deltaE_E_II.png}{0.92}{Preliminary Sample~II MC $\Delta E$--$E$ view retained for layout and analysis review.}{Quarantined analysis output; useful for seeing the intended comparison, not for numerical citation.}

\GatedFigure{figures/gated/fig_b2_b4_two_channel_diagnostic.png}{0.88}{B2--B4 two-channel response diagnostic.}{This is a two-channel correlation diagnostic, not the paper's physical $\Delta E$--$E$ observable.}

\GatedFigure{figures/gated/selected_pulse_inventory.png}{0.88}{Selected-pulse inventory diagnostic.}{The pulse count is a reproducible analysis quantity, not a detector efficiency; claim status remains controlled by the main evidence ledger.}

\subsection{Simulation and method diagnostics}

\PaperFigure{figures/model_diagnostics/timing_mc_method_closure.png}{0.88}{Simulation timing-method closure diagnostic.}{Method-closure simulation only; not a beam-data timing-performance measurement.}

\PaperFigure{figures/model_diagnostics/pid_mc_validation.png}{0.88}{Simulation particle-identification validation diagnostic.}{Simulation result used to understand analysis separation and model behavior; not a beam-data species claim.}

\PaperFigure{figures/model_diagnostics/adc_mc_calibration.png}{0.88}{Simulation ADC-response calibration diagnostic.}{Model diagnostic only; no absolute beam-data calibration is promoted by inclusion here.}

\PaperFigure{figures/model_diagnostics/birks_mc_comparison.png}{0.88}{Birks-quenching simulation comparison.}{Simulation diagnostic illustrating response-model dependence.}

\PaperFigure{figures/model_diagnostics/pileup_digitizer_mc.png}{0.88}{Pile-up and digitizer simulation diagnostic.}{Simulation-only sensitivity view; retained to show the possible supplementary-paper footprint.}

\subsection{Illustrative readout aids}

\PaperFigure{figures/illustrative/03_waveform_annotated.png}{0.88}{Annotated waveform example used to define readout terminology.}{Illustrative only; no quantitative detector claim is attached to the example waveform.}

\PaperFigure{figures/illustrative/06_timewalk_explained.png}{0.88}{Illustration of leading-edge time walk and the motivation for constant-fraction timing.}{Illustrative method aid rather than a measured timing result.}

\subsection{Editorial triage for the next paper pass}

For the main narrative, the preferred direction is to compress redundant figure families rather than promote every diagnostic permanently. In particular, the full-versus-sparse species-resolved $\Delta E$--$E$ set should ultimately become a compact comparison, the beam depth profile should carry one primary topology message plus a concise threshold-sensitivity panel, and the timing figure should foreground the residual distribution and robust interval without turning the median offset into a calibration claim. Until those redesigns are complete, this rough appendix keeps the existing nonblank artifacts visible so collaborators can judge story flow, page count and likely supplementary material.
